% Ad Familiares 10.20 (ad-familiares-10-20)
% Source: https://raw.githubusercontent.com/PerseusDL/canonical-latinLit/master/data/phi0474/phi056/phi0474.phi056.perseus-lat1.xml
% Fetched by scripts/fetch_latin.py

% Dateline (Perseus): Scr. Romae iv K. Iun. a. 711 (43).

\ciceroLetterOpener{CICERO PLANCO}

\ciceroSection{1} ita erant omnia quae istim adferebantur incerta, ut quid ad te scriberem non occurreret. modo enim quae vellemus de Lepido, modo contra nuntiabantur ; de te tamen fama constans nec decipi posse nec vinci ; quorum alterius fortuna partem habet quandam, alterum proprium est prudentiae tuae.

\ciceroSection{2} sed accepi litteras a conlega tuo datas Idibus Maiis, in quibus erat te ad se scripsisse a Lepido non recipi Antonium ; quod erit certius, si tu ad nos idem scripseris. sed minus audes fortasse propter inanem laetitiam litterarum superiorum. verum , ut errare, mi Plance, potuisti (quis enim id effugerit?), sic decipi te non potuisse quis non videt? nunc vero etiam iam erroris causa sublata est ; culpa enim illa 'his ad eundem' vulgari reprehensa proverbio est. sin ut scripsisti ad conlegam, ita se res habet, omni cura liberati sumus nec tamen erimus prius quam ita esse tu nos feceris certiores.

\ciceroSection{3} mea quidem, ut ad te saepius scripsi, haec sententia est, qui reliquias huius belli oppresserit, eum totius belli confectorem fore ; quem te et opto esse et confido futurum. studia mea erga te, quibus certe nulla esse maiora potuerunt, tibi tam grata esse quam ego putavi fore, minime miror vehementerque laetor. quae quidem tu, si recte istic erit, maiora et graviora cognosces. iiii K. Iun.
